% ------------------------------------------------------------------
% Shared preamble for the Metric Spaces & Topology notes.
% Used by notes-ch0-5.tex (full notes) and notes-ch3-5.tex.
% ------------------------------------------------------------------
\usepackage[a4paper,margin=2.5cm,headheight=14pt]{geometry}
\usepackage[T1]{fontenc}
\usepackage[utf8]{inputenc}
\usepackage{lmodern}
\usepackage{microtype}
\usepackage{amsmath,amssymb,amsthm,mathtools}
\usepackage{enumitem}
\usepackage{booktabs,tabularx,array,longtable}
\usepackage[dvipsnames]{xcolor}
\usepackage{tikz}
\usetikzlibrary{arrows.meta,calc,decorations.pathreplacing,patterns}
\usepackage[most]{tcolorbox}
\usepackage{fancyhdr}
\usepackage{titlesec}
\usepackage{amsfonts}
\usepackage{pifont}
\usepackage[hidelinks]{hyperref}
\usepackage{bookmark}

% ---------------- page style ----------------
\pagestyle{fancy}
\fancyhf{}
\fancyhead[L]{\small\itshape Metric Spaces \& Topology}
\fancyhead[R]{\small\itshape \leftmark}
\fancyfoot[C]{\small\thepage}
\renewcommand{\headrulewidth}{0.4pt}
\renewcommand{\chaptermark}[1]{\markboth{\chaptertitlename\ \thechapter.\ #1}{}}
\fancypagestyle{plain}{\fancyhf{}\fancyfoot[C]{\small\thepage}\renewcommand{\headrulewidth}{0pt}}

\titleformat{\chapter}[display]
  {\normalfont\bfseries}{\large\scshape \chaptertitlename\ \thechapter}{4pt}{\LARGE}
  [\vspace{2pt}\titlerule]
\titlespacing*{\chapter}{0pt}{0pt}{18pt}
\titleformat{\section}{\normalfont\large\bfseries}{\thesection}{0.8em}{}
\titleformat{\subsection}{\normalfont\normalsize\bfseries}{\thesubsection}{0.8em}{}

\setlist{itemsep=2pt,topsep=4pt}
\setlength{\parindent}{0pt}
\setlength{\parskip}{5pt}

% ---------------- colours ----------------
\definecolor{defbg}{HTML}{EEF3F9}\definecolor{defrule}{HTML}{2F5D8A}
\definecolor{thmbg}{HTML}{F3EEF9}\definecolor{thmrule}{HTML}{6B3FA0}
\definecolor{exbg}{HTML}{EDF6F0}\definecolor{exrule}{HTML}{2E7D4F}
\definecolor{warnbg}{HTML}{FBEFE8}\definecolor{warnrule}{HTML}{A6461B}
\definecolor{tipbg}{HTML}{FBF6E2}\definecolor{tiprule}{HTML}{8A6D00}

% ---------------- theorem environments ----------------
\theoremstyle{plain}
\newtheorem{theorem}{Theorem}[chapter]
\newtheorem{lemma}[theorem]{Lemma}
\newtheorem{proposition}[theorem]{Proposition}
\newtheorem{corollary}[theorem]{Corollary}
\theoremstyle{definition}
\newtheorem{definition}[theorem]{Definition}
\newtheorem{example}[theorem]{Example}
\newtheorem{exercise}{Exercise}[chapter]
\theoremstyle{remark}
\newtheorem{remark}[theorem]{Remark}
\newtheorem*{warning}{Warning}
\newtheorem*{examtip}{Exam tip}
\newtheorem*{intuition}{Intuition}

\newcommand{\sidebox}[3]{%
  \tcolorboxenvironment{#1}{breakable,enhanced,colback=#2,colframe=#3,
    boxrule=0pt,leftrule=2.5pt,arc=0pt,outer arc=0pt,
    left=6pt,right=6pt,top=3pt,bottom=3pt,before skip=9pt,after skip=9pt}}
\sidebox{definition}{defbg}{defrule}
\sidebox{theorem}{thmbg}{thmrule}
\sidebox{lemma}{thmbg}{thmrule}
\sidebox{proposition}{thmbg}{thmrule}
\sidebox{corollary}{thmbg}{thmrule}
\sidebox{example}{exbg}{exrule}
\sidebox{warning}{warnbg}{warnrule}
\sidebox{examtip}{tipbg}{tiprule}

\newenvironment{solution}{\begin{proof}[Solution]}{\end{proof}}

% Difficulty tags for exercises
\newcommand{\easy}{\textsf{\small\color{exrule}[Easy]}\ }
\newcommand{\medium}{\textsf{\small\color{tiprule}[Medium]}\ }
\newcommand{\hard}{\textsf{\small\color{warnrule}[Hard]}\ }

% Study-plan slot at the start of each chapter
\newtcolorbox{studyslot}[1]{enhanced,colback=white,colframe=defrule!60,
  boxrule=0.6pt,arc=2pt,left=6pt,right=6pt,top=3pt,bottom=3pt,
  fonttitle=\sffamily\bfseries\small,coltitle=defrule,colbacktitle=defbg,
  title={Study slot: #1}}

% ---------------- notation ----------------
\newcommand{\R}{\mathbb{R}}
\newcommand{\N}{\mathbb{N}}
\newcommand{\Q}{\mathbb{Q}}
\newcommand{\Z}{\mathbb{Z}}
\newcommand{\C}{\mathbb{C}}
\newcommand{\eps}{\varepsilon}
\newcommand{\cl}{\operatorname{cl}}
\newcommand{\inte}{\operatorname{int}}
\newcommand{\ext}{\operatorname{ext}}
\newcommand{\diam}{\operatorname{diam}}
\newcommand{\Bo}{B}
\newcommand{\Bc}{\overline{B}}
\newcommand{\abs}[1]{\left\lvert #1\right\rvert}
\newcommand{\norm}[1]{\left\lVert #1\right\rVert}
\newcommand{\set}[2]{\left\{\, #1 \;:\; #2 \,\right\}}
\newcommand{\seq}[1]{(#1_n)_{n\in\N}}
\newcommand{\dinf}{d_\infty}
\newcommand{\Ccal}{\mathcal{C}}
\newcommand{\U}{\mathcal{U}}
\newcommand{\topo}{\tau}
\newcommand{\ip}[2]{\left\langle #1, #2\right\rangle}
\newcommand{\cmark}{\ding{51}}
\newcommand{\xmark}{\ding{55}}
\newcommand{\tickbox}{$\square$\ }
\newcolumntype{Y}{>{\raggedright\arraybackslash}X}
